% Eigenständig kompilierbare Aufgabenstellung.
% Kompilieren mit XeLaTeX oder LuaLaTeX (wegen fontspec/polyglossia), z. B.:
%   latexmk -xelatex aufgabenstellung_standalone.tex
\documentclass[a4paper,11pt]{article}

\usepackage{polyglossia}
\setdefaultlanguage[spelling=new,babelshorthands=true]{german}
\usepackage{fontspec}
\usepackage{parskip}
\usepackage[left=2.5cm,right=2.5cm,top=2.5cm,bottom=2.5cm]{geometry}

\title{Aufgabenstellung}
\author{}
\date{}

\begin{document}

\begin{center}
  {\Large\bfseries Entwicklung eines funkbasierten Beleuchtungssystems\\
  auf Basis von LED-Röhren für die Veranstaltungstechnik}\\[0.5em]
  {\large Aufgabenstellung}
\end{center}

\bigskip

Der Aufbau professioneller Bühnenbeleuchtung ist durch die kabelgebundene
DMX-Signalverteilung mit einem erheblichen Planungs-, Verkabelungs- und
Zeitaufwand verbunden. Kommerzielle Funklösungen adressieren dieses Problem,
sind jedoch kostenintensiv und primär auf große Produktionen ausgerichtet.
Ziel dieser Arbeit ist die Entwicklung eines kostengünstigen, funkbasierten
Beleuchtungssystems auf Basis von LED-Leuchtröhren, das DMX-Steuersignale
drahtlos an mehrere Beleuchtungseinheiten überträgt und sich ohne Anpassungen
in bestehende DMX-Infrastrukturen integrieren lässt. Das System soll dabei
durchgängig als frei zugängliches Open-Source-Projekt konzipiert und
veröffentlicht werden, um einen niedrigschwelligen Nachbau zu ermöglichen.
Das System besteht aus einer zentralen \emph{Bridge}, welche Steuerdaten
entgegennimmt und verteilt, sowie den empfangenden LED-Leuchtröhren, die die
Daten unmittelbar in Lichtsignale umsetzen.

Im Rahmen der Arbeit sind folgende Teilaufgaben zu bearbeiten:

\begin{itemize}
  \item Analyse des Einsatzszenarios und Ableitung funktionaler sowie
        nicht-funktionaler Anforderungen (u.\,a. Echtzeitverhalten,
        Zuverlässigkeit, Reichweite, Benutzerfreundlichkeit).
  \item Entwurf einer Systemarchitektur aus Bridge und Leuchtröhren auf Basis
        des ESP32 mit Definition eines schlanken, drahtlosen
        Übertragungsprotokolls (ESP-NOW) für die DMX-Datenverteilung.
  \item Implementierung der Firmware für die Betriebsmodi Einzelbetrieb,
        autonomer Master-Modus sowie drahtloser DMX-Modus mit Empfang
        kabelgebundener DMX512-Signale und Weiterleitung an adressierbare Leuchtröhren.
  \item Umsetzung der Ansteuerung der WS2815-LED-Strips, einer menügeführten
        Bedienung über OLED-Display und Taster sowie einer persistenten
        Konfigurationsspeicherung.
  \item Entwurf der Hardware in Form eigener Platinenlayouts (PCB) sowie der
        3D-Gehäuse für Bridge und Leuchtröhren unter Berücksichtigung der
        elektrischen Belastbarkeit und Portabilität.
  \item Experimentelle Evaluierung des Systems in einer realitätsnahen
        Umgebung hinsichtlich Verbindungsstabilität, Übertragungsrate, Latenz
        und Jitter sowie Ableitung verbleibender Optimierungspotenziale.
\end{itemize}

Sämtliche Ergebnisse -- Software, Platinenlayouts und Gehäusemodelle -- werden
als Open-Source-Projekt veröffentlicht, um Nachvollziehbarkeit und
Weiterentwicklung durch Dritte zu ermöglichen.

\end{document}
